% !TeX root = ../../main.tex
\subsection{Stable trees and their tails}
\label{subsec:s06:03}

For a vertex $v$ of a tree, put $m_v=2g_v-3+n_v$. Let
$\mathbf s_v$ be its tuple of flag parameters: both flags belonging
to an edge $e$ use the same $s_e\in[0,1]$.
For the comparisons in this and the next subsection define the
positive scale
\begin{equation}\label{eq:s06:tree-scale}
 \mathcal T_g=\beta_g\frac{g_2!R^{-g_2}}{\sqrt{g_2+1}}.
\end{equation}
We first keep the exact automorphism factors. Only after bounding
every vertex will we replace unrooted trees by plane rooted trees
in a positive majorant. This controls the entire genus and valence
sum at each edge count.

\begin{prp}[The complete tree sum]\label{prp:s06:tree-tail}
For $g\geq2$,
\begin{equation}\label{eq:s06:tree-expansion}
 \frac{12^{g_2}b_g}{g_2!}
 =\sum_{\substack{\Gamma\in\mathcal G_{g,0}\\\Gamma\text{ a tree}}}
  \frac1{|\operatorname{Aut}\Gamma|}
  \int_{[0,1]^{E(\Gamma)}}
       \prod_{v\in V(\Gamma)}W^{\rm top}_{g_v,n_v}(\mathbf s_v)
       \prod_{e\in E(\Gamma)}\dd s_e.
\end{equation}
Every contributing tree has $E\leq g_2$ edges. There is a constant
$C>0$, independent of $g$ and $E$, such that the absolute sum of
its $E$-edge integrals is at most
\begin{equation}\label{eq:s06:edge-mass}
 \mathcal T_g C^{E+1}
 \sqrt{\frac{g_2+1}{g_2-E+1}}\,
                  \frac{(g_2-E)!}{g_2!}
 \qquad(0\leq E\leq g_2).
\end{equation}
Hence, for every fixed integer $K\geq1$, the sum over all trees
with at least $K$ edges is $O_K(\mathcal T_g g_2^{-K})$ in
absolute value.
\end{prp}
\begin{prf}
\emph{The finite-genus identity.}
The boundary exponential with both scalar coefficients one gives
edge kernel $\int_0^1e^{-s(\psi_f+\psi_{f'})}\dd s$.
By Corollary~\ref{crl:s02:defect-allocation}, the pullback of
$\lambda_g$ vanishes on every graph with a cycle and is the product
of the vertex top classes on a tree. Logarithmic kappa additivity
\eqref{eq:s02:tautological-gluing} now gives the right side of
\eqref{eq:s06:tree-expansion}, with the original graph
$1/|\operatorname{Aut}\Gamma|$.
On the left, $\kappa_1+\delta=12\lambda_1$ and degree conversion
gives $\int\lambda_ge^{12\lambda_1}=12^{g_2}b_g/g_2!$.
For a tree, the genus condition and flag count imply
\begin{equation}\label{eq:s06:tree-degree}
 \sum_v m_v=\sum_v(2g_v-3+n_v)=g_2-E.
\end{equation}
All $m_v$ are nonnegative, proving $E\leq g_2$.

\emph{A majorant for every tree and genus assignment.}
The comparison for $\beta_G$ in Theorem~\ref{thm:s06:lambda-g}
uses $\sum_v g_v=g$ to give
$\prod_v\beta_{g_v}\leq C^{E+1}\beta_g$, including rational and
elliptic vertices. Apply Lemma~\ref{lmm:s06:vertex-bound} to
every factor. Since $\sum_v n_v=2E$ and there are $E+1$ vertices,
all its exponential constants are bounded by $C^{E+1}$ after
increasing $C$. The edge domains have measure one. After removing
$\beta_gR^{-g_2}$, the remaining majorant is thus a constant to
power $E+1$ times
$\prod_v m_v!/\sqrt{m_v+1}$.

Every unrooted genus-labeled tree has a plane rooted realization,
and its original weight is at most one. A plane rooted tree with
$E$ edges is encoded by a walk of $2E$ up/down steps, so there are
at most $4^E$ shapes. For a fixed shape and the $m_v$, the genus
labels are determined by $2g_v=m_v+3-n_v$. The beta factors have
already been bounded uniformly; dropping parity and stability restrictions
therefore only enlarges the positive sum over nonnegative $m_v$.
This introduces no factorial in the vertex count.

For every integer $M\geq0$,
\begin{equation}\label{eq:s06:factorial-convolution}
 \sum_{k=0}^M k!(M-k)!\leq3M!.
\end{equation}
Indeed the two endpoint terms give $2M!$ when $M\geq1$, while
$\binom{M}{k}\geq M$ for $1\leq k\leq M-1$; the case $M=0$
is immediate.

Iterating the convolution for $E+1$ vertices gives
$3^EM!$ for total $M=\sum_v m_v$. Also
$\prod_v(m_v+1)\geq M+1$. With $M=g_2-E$, these estimates and
the shape count prove \eqref{eq:s06:edge-mass}.

\emph{Summing every edge count.}
For $0\leq E\leq g_2$, direct multiplication gives
\[
 (E+1)(g_2-E+1)-(g_2+1)=E(g_2-E)\geq0.
\]
The square root in \eqref{eq:s06:edge-mass} is therefore at most
$\sqrt{E+1}$. Moreover
\[
 \frac{(g_2-E)!}{g_2!}\leq\left(\frac e{g_2}\right)^E.
\]
For example, integrate $\log t$ from $g_2-E$ to $g_2$ to bound
the logarithm of the falling factorial below by
$E(\log g_2-1)$. For sufficiently large $g$, $Ce/g_2\leq1/2$.
Thus the sum of \eqref{eq:s06:edge-mass} over $E\geq K$ is at most
\[
 C\mathcal T_g\left(\frac{Ce}{g_2}\right)^K
       \sum_{j\geq0}2^{-j}\sqrt{K+j+1}
 =O_K(\mathcal T_g g_2^{-K}).
\]
This bounds the complete tree sum, including all small-genus and
large-valence configurations.
\end{prf}

A one-edge tree has positive genus at each endpoint. Its two
orientations give the exact expression
\begin{equation}\label{eq:s06:one-edge}
 \frac12\sum_{G=1}^{g-1}\int_0^1
       W^{\rm top}_{G,1}(s)W^{\rm top}_{g-G,1}(s)\dd s.
\end{equation}
The factor $1/2$ also gives the correct automorphism weight when the
two genera agree. The next bound isolates the elliptic endpoints
without a fixed-split approximation.

\begin{lmm}[Nonelliptic one-edge tail]\label{lmm:s06:nonelliptic-tail}
The absolute contribution to \eqref{eq:s06:one-edge} of all splits
with $2\leq G\leq g-2$ is $O(\mathcal T_g g_2^{-3})$.
\end{lmm}
\begin{prf}
At an endpoint of genus $G\geq2$, the coefficient order is
$2G-2\geq2$. The two orders sum to $M=g_2-1$.
Applying the all-type vertex and beta bounds as above, using
$(m+1)(M-m+1)\geq M+1$, bounds the relative contribution by
\[
 C\sqrt{\frac{g_2+1}{M+1}}\,
       \frac{M!}{g_2!}
       \sum_{m=2}^{M-2}\binom{M}{m}^{-1}.
\]
Allowing odd $m$ only enlarges the sum. For $M\geq6$ it is at most
\[
 \frac2{\binom M2}+\frac{M-5}{\binom M3}=O(M^{-2}).
\]
Since $M!/g_2!=1/g_2$, the asserted bound follows. The finitely
many smaller totals are absorbed in the constant.
\end{prf}

\begin{exm}[Genus two]
Here $g_2=1$. The tree with no edges has value $\beta_2$, and the only
one-edge tree consists of two elliptic vertices with weight $1/2$.
The series \eqref{eq:s06:beta-series} gives
$\beta_1=1/24$ and $\beta_2=7/5760$. Thus
\[
 b_2=\frac{\beta_2+\beta_1^2/2}{12}=\frac1{5760}.
\]
This checks both the stack symmetry factor and the conversion by
$12^{g_2}$.
\end{exm}
